%%%%%%%%%%%%%%%%%%%%%%%%%%%%%%%%%%%%%%%%%%%%%%%%%%%%%%%%%%%%%%%%%%
% proposal-rebecca.tex -- exploratory PhD proposal for discussion with
% Dr. Rebecca Caines (York/TMU Communication & Culture), revised after the
% October 2026 review (see CHANGELOG-review-2026-10.md).
%
% Layout follows the York thesis template supplied by the author
% (letter paper, 11 pt, 25 mm margins, one-and-a-half spacing, no
% paragraph indent, Times-style text, no institutional logo).
% Deviations from that template: author-year citations (natbib/plainnat)
% instead of numeric ieeetr, and a small paragraph skip because the
% template sets \parindent to 0.
% Compiles with pdfLaTeX, XeLaTeX or LuaLaTeX; bibliography via BibTeX.
% main.tex simply \input{}s this file so Overleaf compiles it by default.
%%%%%%%%%%%%%%%%%%%%%%%%%%%%%%%%%%%%%%%%%%%%%%%%%%%%%%%%%%%%%%%%%%
\documentclass[letterpaper,11pt]{article}
\usepackage{iftex}
\ifPDFTeX
  \usepackage[T1]{fontenc}
  \usepackage[utf8]{inputenc}
  \usepackage{mathptmx}
\else
  \usepackage{fontspec}
  \setmainfont{texgyretermes}[Extension=.otf, UprightFont=*-regular,
    BoldFont=*-bold, ItalicFont=*-italic, BoldItalicFont=*-bolditalic]
\fi
\usepackage[top=25mm, bottom=25mm, left=25mm, right=25mm]{geometry}
\usepackage{setspace}
\usepackage{enumitem}
\usepackage{booktabs}
\usepackage{tabularx}
\usepackage{array}
\usepackage[round,authoryear]{natbib}
\usepackage{xurl}
\usepackage[hidelinks]{hyperref}

\setlength{\parindent}{0pt}
\setlength{\parskip}{6pt}
\setlength{\emergencystretch}{1.5em}
\renewcommand{\refname}{Bibliography}
\setlist{itemsep=2pt, topsep=4pt, parsep=0pt}
\onehalfspacing
\widowpenalty=10000
\clubpenalty=10000

\begin{document}
\begin{center}
  {\bf\Large Whose Contribution Counts?}\\[4pt]
  {\bf\large AI Mediation and Collective Authorship in Digital Improvisation}\\[14pt]
  {\bf Roman Jahandideh}\\[4pt]
  Exploratory proposal for discussion with Dr.\ Rebecca Caines\\
  Graduate Program in Communication \& Culture, York University\\
  Intended Fall 2027 entry\\[4pt]
  \today
\end{center}

\begin{quote}
\small\textit{Working proposal prepared for discussion with Dr.\ Rebecca Caines.
The research setting and project partnerships will be developed in
consultation with the prospective supervisor.}
\end{quote}

\section{Problem and cultural stakes}

Generative models now sit inside tools that people use to make things
together, and in these tools a model does more than produce content: it
\emph{reinterprets} what people contribute. Someone types a phrase, and a
moment later it has become a sound or a moving form in a shared scene. The
group then has to decide what that transformation means, and whether it
still belongs to the person who offered it.

The stakes reach beyond any single workshop. When AI translates creative
contributions into shared audiovisual material, its outputs may become more
visible or influential than the contributions that prompted them. I treat
this as a possibility to investigate, not an established effect. If it does
happen, it raises a cultural question about whose interpretations gain
authority, whose contributions remain recognisable, and how disagreement is
negotiated in collective creative practice.

Improvisation is a good place to ask this question, because improvisers
already have ways of listening to, building on and refusing one another's
offers. Improvisation has also been argued to be a research method in its own
right, one grounded in risk, active listening and collaborative response
\citep{caines2020fragile}. A machine that transforms offers enters these
practices, and the people involved have to work out where it fits.

\section{Research gap}

Existing research examines human--AI co-creation, collective decision-making
and the attribution of creative credit. This project investigates their
intersection in small-group digital improvisation: how a contribution becomes
recognisable as someone's work, how its transformation can be challenged, and
how authority over the resulting artifact is negotiated over repeated
sessions. Three studies come closest to this question, and each leaves a
different part of it open.

\citet{xu2026} observed six groups of three university students designing a
cultural educational game around one shared ChatGPT interface, in a single
45-minute session, and interviewed nine further students. Shared AI output
became an anchor for attention, a shared memory and a negotiable part of the
groups' decisions. The authors note that their analysis emphasised
information flow over social aspects such as trust and accountability. I take
their question into creative work in which a specific person's material is
transformed, and I follow authorship and accountability across repeated
sessions.

\citet{trajkova2024} ran focus groups with 24 dance students about partnered
movement improvisation, in order to design an AI dance partner. In that work
the AI is a \emph{partner}. In mine it is a \emph{mediator}, reinterpreting
one person's contribution for a whole group.

\citet{epstein2020} showed, through online vignette experiments, that
describing an AI as an agent rather than a tool shifts the credit and
responsibility that observers assign to the people behind an AI artwork.
Their participants judged a finished work from the outside. I ask how credit
and responsibility are claimed and contested \emph{in situ}, by the people
making the work.

Questions of contestation point the same way. Contestability has mainly been
developed for high-stakes algorithmic decisions such as hiring or sentencing
\citep{lyons2021}. Much less has been said about what it means to challenge a
machine's \emph{creative} reading of one's contribution, in front of
collaborators.

\section{Research questions}

\textbf{Central question.} In collaborative digital improvisation, how do
participants, facilitators and a generative-AI interface shape whether a
person's contribution is recognised, transformed, taken up, contested and
credited in the shared work?

\begin{enumerate}[label=\textbf{RQ\arabic*.}, leftmargin=*]
  \item How do participants interpret differences between what they offered
        and how it was transformed?
  \item How do groups and facilitators negotiate acceptance, revision,
        disagreement and refusal, and who comes to have the standing to
        decide?
  \item How do these negotiations shape felt ownership, public attribution
        and creative responsibility for the resulting work?
\end{enumerate}

Two terms need defining. On \textbf{participation}: I focus on the trajectory
of identifiable creative contributions, while also attending to listening,
facilitation, silence and refusal as forms of participation. Recognition will
be examined through participants' interpretations and interactions, rather
than inferred solely from incorporation into the artifact. I do not assume
that group decision-making is fair, that individual control is enough, or
that facilitation does harm. A majority can override someone; a facilitator
can help a quieter participant take part.

By \textbf{authorship} I mean three related things: \emph{felt ownership},
whether a person experiences part of the work as theirs; \emph{public
attribution}, how the work is credited; and \emph{creative responsibility},
who answers for what the work does. Legal ownership and copyright are outside
the scope of this study.

\section{Theoretical framework}

My primary framework is Lucy Suchman's study of human--machine
reconfiguration \citep{suchman2007}, which examines how agencies are figured
at the interface and how humans and machines come to be treated as similar or
different in practice. Read this way, the question of whether a transformed
element ``belongs'' to its contributor, to the group or to the system is not
settled in advance. People work it out as they interact. This gives me three
questions to bring to the material:
\begin{itemize}
  \item how participants attribute agency and authorship in talk and action
        (``the AI did that'', ``that one is mine'');
  \item how interface rules, model defaults, facilitation and my own design
        decisions make some actions easy and others difficult;
  \item who is held responsible when a transformation misrepresents someone.
\end{itemize}
Improvisation studies adds a limited, complementary lens: attention to
listening and response, and to how an offer is taken up, ignored or refused
\citep{caines2020fragile}. These concepts tell me where to look. They do not
decide what I will find.

\section{Relationship to my master's research}

At Simon Fraser University, my master's research has asked how interfaces
make a person's state perceptible to others. In a cooperative Unity game, I
rendered remote teammates' heart rate and heart-rate variability as
label-free ambient icons; in a second project, I built a Unity environment
that responds to the player's heart rate. Both manuscripts were submitted to
ACM venues and received revision requests. Neither has been accepted
\citep{jahandideh_teammate_ms,jahandideh_bioterrarium_ms}.

The teammate study left me with a question I could not answer inside it.
Participants tied their willingness to use the cue to control over it: they
asked for opt-in sharing, an immediate off switch and separate permissions
for the icon and the numbers \citep{jahandideh_teammate_ms}. Looking back as
the designer, I also noticed that the classification behind the label-free
display still sorted people into evaluative categories they never saw. Both
observations lead to the question this PhD takes up: who controls how a
shared system interprets people?

The PhD does not continue physiological sensing. It moves from involuntary
bodily signals to deliberate creative contributions, and from consent to
being made visible to the negotiation of being transformed. My training fits
this shift. In architecture and landscape architecture I learned to design
how people perceive and move through space; in HCI I have built and evaluated
networked interactive systems myself. That combination prepares me to design
a shared environment as a creative work, and then to study how people
interpret it.

\section{Research-creation methodology}

The project follows the program's research-creation option, which combines a
creative work with a written dissertation of 150--200 pages that develops the
argument \citep{comcult_rc}. Of the modes \citeauthor{chapman2012}
\citeyearpar{chapman2012} describe, it is closest to
``research-from-creation'': making the work generates material for
understanding how participants engage with it.

The creative component will comprise the shared audiovisual environment and
the improvised works produced through it. Developing and revising its
expressive possibilities, interaction rules and workshop scores will form
part of the inquiry into how creative authority is distributed. What
participants make are short collective audiovisual improvisations: projected,
spatial scenes of light, form and sound that they build and perform together
in real time, and then talk about.

A brief, illustrative session shows what this looks like. Four improvisers
receive a score: build the weather of a city none of them has visited. One
types ``the bus that never comes.'' The system proposes a slow amber pulse at
the edge of the space with a low drone, and displays its generated
interpretation, ``waiting as a distant repeating light.'' The contributor
accepts it. Later, another participant wants the pulse brighter. The
contributor declines, and the facilitator asks the group how to proceed.
Moments like these, of acceptance, challenge and resolution, are what the
study examines.

The work unfolds in four phases:
\begin{enumerate}[label=\arabic*., leftmargin=*]
  \item \emph{Scoping interviews} with improvisers and facilitators about how
        offers, refusal and credit work in their practice, with and without
        technology. What I learn shapes the scores and the system's rules.
  \item \emph{Co-design} sessions with a few practitioners to choose the input
        form, set limits on the transformations, design controls for refusal
        and credit, and test accessibility provisions.
  \item \emph{Repeated workshops} with two or three small groups, each meeting
        several times so that shared norms can form. The number of groups and
        sessions will depend on access and on the material that emerges.
  \item \emph{Reflection interviews} in which participants review video
        episodes from their own sessions.
\end{enumerate}

Each group works under three arrangements, in an order that varies across
sessions:

\begin{center}
\small
\begin{tabularx}{0.92\linewidth}{@{}>{\raggedright\arraybackslash}p{3.3cm}>{\raggedright\arraybackslash}X>{\raggedright\arraybackslash}X@{}}
\toprule
Arrangement & Who proposes the transformation & Who approves it \\
\midrule
Contributor-led & AI system & The contributor \\
Group-led & AI system & The group, by a procedure it agrees in its first session \\
Human-mediated & Another participant in a rotating ``interpreter'' role, using the same asset library & The contributor (same rule as contributor-led) \\
\bottomrule
\end{tabularx}
\end{center}

The contributor-led and group-led arrangements explore different
distributions of approval authority. The human-mediated arrangement, which
uses comparable assets and the same approval rule as the contributor-led one,
offers an additional interpretive contrast. It also changes timing, social
relationships and expectations. These workshops will therefore support
situated comparison rather than isolate a causal effect of AI, and a study of
this size does not generalise statistically.

\section{The system and the place of AI}

A participant submits a short text phrase; sketches may be added later if
co-design supports them. An existing language model returns one
\emph{proposal}, a choice of elements and settings from a curated library of
forms, light, colour, particle ``weather'', sound and text, together with a
one-line \emph{generated interpretation} of the phrase that everyone can
see. That line is part of the model's output. It does not explain why the
model produced the transformation, which it cannot establish. Depending on
the arrangement, contributors or groups can inspect, edit, accept, reject,
request another proposal or undo. A manual path to the same library is always
open, so no one needs AI in order to contribute.

Generative AI is included because machine-mediated reinterpretation is the
object of study. The project does not presume that AI improves improvisation;
it examines how participants respond to its transformations and negotiate
their place in the shared work. The model does not decide whether anything
enters the work, who is credited or what is removed, and it makes no
inferences about participants' emotions or behaviour.

\textbf{Feasibility.} Unity applies only a fixed set of permitted actions
with bounded settings. Model output is checked against that set and never run
as code; malformed output or a service failure simply produces no proposal.
Each proposal and decision is recorded and can be undone. I will first test
a locally run open-weight model, which avoids third-party processing, and
otherwise use a hosted service under institutionally approved terms, sending
only the phrase. I will not train or fine-tune models: this is a
research-creation project built on existing tools, not an AI engineering
dissertation.

\section{Case-study context and access}

\textbf{Possible primary context, subject to Dr.\ Caines's agreement:}
workshops with adult improvising artists and creative-technology
practitioners in Toronto, connected to her SSHRC Insight project
\emph{Interfacing, Playtesting, Resetting: Reframing Technology Use Through
Improvised Creative Technologies} \citep{yfile2026sshrc}. Only the project's
title and funding are public. I do not assume that it includes this study,
these participants or any AI component; whether a student study could sit
alongside it is for her to judge.

\textbf{Alternative:} open calls to Toronto improvisation, theatre and sound
communities and to AMPD students, excluding anyone I teach, with an arts
organisation as host only if it agrees. Earlier projects such as multiPLAY
are precedents for digitally mediated improvisation, not recruitment sources.
I will not work with disability-focused or Indigenous communities without
established partnerships, community governance and relevant expertise.

\section{Analysis}

I will construct multimodal accounts of selected episodes, linking
participants' words and actions with screen events and system records. An
episode runs from a contribution, through its transformation and the group's
response, to the next creative action. Analysis will compare accepted
transformations, contested transformations, refusals, and episodes in which
participants disagree about authorship. Later interviews will add
participants' own interpretations, including accounts that challenge my
initial reading, and comparisons across sessions will show how these
practices change as groups establish shared norms.

For the episodes I adapt the video microethnography that \citet{mcarthur2015}
used to study players' interactions with interfaces. Interviews will be
analysed through reflexive thematic analysis \citep{braun2019}. System
records show what happened in the software; they are not a ``true'' account
of intention against which participants are judged. When records,
observation and participants' accounts diverge, the divergence is itself a
finding.

\section{Ethics, consent and accessibility}

Consent will be specific, ongoing and revisable
\citep{im2021,strengers2021}. Participants consent separately to recording,
to AI processing of their contributions, and to public dissemination of the
works.

Withdrawal needs particular care in collective work. Provenance will support
identifying and, where feasible, removing attributable contributions.
Participants will agree on withdrawal procedures and deadlines before the
workshops, including the limits of removing contributions from jointly
developed or publicly disseminated works. Credit for public works will be
agreed with participants, and their time will be paid. Data will be
minimised and pseudonymised, any third-party processing disclosed, and
retention periods stated.

Full participation without AI is a design commitment that the workshops will
test. Planned provisions are a reduced-motion setting, readable and resizable
text, visual or textual alternatives to sound-only information, controls that
a facilitator can operate on a participant's behalf, and the manual path.

Image-generation models can reproduce stereotypes even when prompts avoid
demographic terms \citep{bianchi2023}, so I will discuss model defaults
openly with participants. Participation alone does not make a project
equitable \citep{birhane2022}. Finally, as designer, researcher and possibly
facilitator, I shape the situation I study. I will keep a design journal
and, where possible, work with a co-facilitator.

\section{Intended contributions}

The intended scholarly contribution is an account of how creative authority
is negotiated across contributors, facilitators, interface rules and
AI-generated transformations. The study will examine when technical
attribution supports recognition, and when it fails to resolve disagreement
about authorship. The creative system and workshop documentation will make
these tensions available for analysis and further practice.

These are aims, not findings. The workshops may show that AI is unnecessary
in some arrangements. If so, the contribution would lie in explaining when,
how and why it becomes unnecessary, and what that reveals about collective
creativity.

\section{Provisional chapter structure}

\begin{enumerate}[label=\arabic*., leftmargin=*]
  \item Introduction.
  \item Co-creation, credit and contestation: literature.
  \item Framework: human--machine reconfiguration and improvised response.
  \item Methodology, research-creation and reflexivity.
  \item The environment as a creative work and a configured situation.
  \item Recognition and transformation.
  \item Disagreement, refusal and facilitation.
  \item Ownership, attribution and responsibility.
  \item Conclusion.
\end{enumerate}

\section{Timeline (full-time, typical program schedule)}

\begin{center}
\small
\begin{tabularx}{\linewidth}{@{}lX@{}}
\toprule
Year 1 (2027--28) & Courses (CMCT 7000, 7200, 7005 and electives); literature review; ethics preparation; technical tests with a local model. \\
Year 2 (2028--29) & Comprehensive examinations; formal dissertation proposal (maximum 3,500 words) with the supervisory committee \citep{york_dissregs}; ethics approval; scoping interviews. \\
Year 3 (2029--30) & Co-design; repeated workshops; reflection interviews; ongoing analysis. \\
Year 4 (2030--31) & Analysis; documentation or public sharing of the creative work; writing and defence. Year 5 is held in reserve. \\
\bottomrule
\end{tabularx}
\end{center}

\section{Program and supervisor fit}

The project belongs to the program's Technology in Practice area, with Media
\& Culture as a minor. It approaches AI-mediated creative work as a cultural
practice of participation, recognition and authorship, rather than as a
technical problem.

Dr.\ Caines's research on improvisation as a research methodology
\citep{caines2020fragile} and on socially responsive creative technologies
\citep{caines2023consultation} speaks directly to the project's setting and
methods. Her recent research-creation study of generative AI in
interdisciplinary classrooms attends to ethics, access, collaboration and
improvisatory skills \citep{caines2026steaming}. My project carries that
concern from learning with these tools to contribution-level recognition and
contested authorship in repeated group improvisation.

\bibliographystyle{plainnat}
\bibliography{references}

\end{document}
